\section{Deep Learning aplicado a audio}\label{sec:dl}

%\subsection{Contextos matemáticos y preprocesamiento}
\subsection{Contextos matemáticos y preprocesamiento}
\label{sec:dl_preproc}

El análisis de señales acústicas mediante \emph{Deep Learning} requiere transformar el audio crudo en representaciones que capturen adecuadamente sus propiedades espectrales y temporales.  
Para ello se emplean herramientas de la teoría de Fourier y representaciones perceptuales que aproximan el comportamiento del sistema auditivo humano.  
En lo que sigue se presentan las bases matemáticas de tres transformaciones fundamentales: la Transformada Rápida de Fourier (FFT), la Transformada de Fourier de Tiempo Corto (STFT) y la escala Mel.

\subsubsection{Transformada Rápida de Fourier (FFT)}

Dada una señal discreta $x[n]$ de longitud $N$, su Transformada Discreta de Fourier (DFT) está definida como
\begin{equation}
	X[k] = \sum_{n=0}^{N-1} x[n] \, e^{-j \tfrac{2\pi}{N} kn}, 
	\quad k = 0,1,\dots,N-1.
	\label{eq:dft}
\end{equation}

La DFT descompone la señal en senoidales complejas ortogonales, revelando la distribución de energía en frecuencia \cite{Haykin2001,Proakis2007}.  
Su cálculo directo tiene costo $\mathcal{O}(N^2)$, pero el algoritmo FFT lo reduce a $\mathcal{O}(N \log N)$, habilitando aplicaciones en tiempo real como detección acústica.  

Es importante notar que la DFT brinda una descripción global: indica magnitud y fase de cada componente frecuencial, pero pierde información sobre el instante en que ocurrió.  
En bioacústica esto es relevante: el zumbido de un mosquito puede variar en fracciones de segundo, lo que motiva el uso de representaciones con resolución conjunta en tiempo y frecuencia.

\insertarfigura{fig/fft_signal.png}
{Conversión de una señal temporal a su representación espectral mediante FFT.}
{fig:fft_signal}{1}

\subsubsection{Transformada de Fourier de Tiempo Corto (STFT)}

La STFT extiende la DFT incorporando localización temporal.  
La señal se multiplica por una ventana $w[n]$ de duración finita, que selecciona un segmento, y se desplaza en el tiempo mediante un parámetro $m$. Su definición es:
\begin{equation}
	X(m,\omega) = \sum_{n=-\infty}^{\infty} x[n] \, w[n-m] \, e^{-j \omega n}.
	\label{eq:stft}
\end{equation}

El resultado es una función bidimensional que describe cómo varía el contenido espectral en el tiempo, dando lugar al \emph{espectrograma}.  
La resolución temporal y frecuencial dependen inversamente de la longitud de la ventana:  
ventanas cortas capturan sucesos breves pero con poca resolución espectral, mientras que ventanas largas ofrecen buena resolución en frecuencia pero pierden detalle temporal.  
Esta dualidad, conocida como principio de incertidumbre tiempo–frecuencia, condiciona el diseño del preprocesamiento \cite{Bancroft2002}.  

En detección de mosquitos, la STFT permite observar la evolución dinámica del zumbido, que puede variar en frecuencia fundamental por el movimiento del insecto o por interferencias ambientales.

\insertarfigura{fig/stft_spectrogram.png}
{Esquema de la STFT mostrando segmentación por ventanas.}
{fig:stft_spectrogram}{1}

\subsubsection{Escala Mel}

Aunque la STFT proporciona un espectro en Hz, el sistema auditivo humano no percibe las frecuencias de manera lineal.  
Somos más sensibles a bajas frecuencias: diferenciar entre 100 y 200 Hz resulta más fácil que entre 5000 y 5100 Hz.  
Para aproximar esta percepción se usa la escala Mel, que relaciona la frecuencia física $f$ con una perceptual $f_{\text{mel}}$:
\begin{equation}
	f_{\text{mel}} = 2595 \, \log_{10}\!\left( 1 + \frac{f}{700} \right).
	\label{eq:mel}
\end{equation}

El cálculo de un \emph{mel-spectrograma} consiste en aplicar un banco de filtros triangulares, distribuidos equidistantemente en la escala Mel, sobre el espectrograma lineal.  
La energía espectral se proyecta a un dominio perceptual comprimido, reduciendo dimensionalidad e introduciendo cierta invariancia frente a variaciones irrelevantes para el oído humano.  

En \emph{Deep Learning} aplicado a audio, el mel-spectrograma es una de las representaciones más utilizadas como entrada a CNNs, pues concentra la información discriminativa y mejora la robustez frente a ruido ambiental \cite{Campos2022,HumBugDB}.

\insertarfigura{fig/mel_spectrogram.png}
{Waveform, espectrograma y mel-spectrograma de un segmento de audio de 10 segundos (Google AudioSet). 
	El mel-spectrograma, basado en la escala perceptual Mel, ofrece mayor resolución en bajas frecuencias que el espectrograma lineal.}
{fig:mel_spectrogram}{0.9}

% Archivo: cap_dl_1.tex
% Contenido: Explicación de FFT, STFT y escala Mel en audio. 
% Relevancia para extracción de características y robustez frente al ruido.
% Referencias: Haykin (2001), Proakis (2007), Bancroft (2002).



\subsection{CNNs aplicadas a HumBugDB}
\label{sec:dl_cnn}

El conjunto de datos \emph{HumBugDB} constituye un recurso extenso y representativo de grabaciones acústicas de mosquitos en entornos de laboratorio y de campo \cite{HumBugDB}.  
La señal característica del mosquito es un zumbido cuasi-periódico producido por el batido de alas, con frecuencia fundamental y armónicos.  
Su detección en grabaciones reales es desafiante, ya que suele estar enmascarada por ruidos ambientales como viento, tráfico o voces humanas.

\subsubsection*{Inferencia bayesiana mediante MC Dropout}

En redes profundas, las predicciones suelen ser tratadas como determinísticas, sin cuantificar la confianza.  
Gal y Ghahramani (2016) demostraron que el \emph{dropout} —originalmente introducido como técnica de regularización— puede interpretarse como una aproximación bayesiana cuando se mantiene activo también durante la inferencia.  

Este procedimiento, conocido como \emph{Monte Carlo Dropout} (MC Dropout), consiste en realizar múltiples pasadas de inferencia con distintas máscaras de dropout, generando una distribución de salidas en lugar de un valor único.  
A partir de esa distribución se obtienen medidas de incertidumbre, como la varianza predictiva, lo que permite distinguir entre predicciones seguras y dudosas \cite{GalGhahramani2016}.  
En audio bioacústico esta capacidad resulta esencial, pues las señales suelen estar contaminadas por ruido.

\subsubsection{Motivación del uso de CNNs}

Las Redes Neuronales Convolucionales (CNNs) son especialmente adecuadas porque trabajan sobre representaciones bidimensionales como los mel-spectrogramas.  
Una capa convolucional aplica filtros $w^{(k)}$ sobre la entrada $x$ de acuerdo con\(
y_{i,j}^{(k)} = \sigma \!\left( \sum_{m=0}^{M-1}\sum_{n=0}^{N-1} 
w_{m,n}^{(k)} \, x_{i+m,\,j+n} + b^{(k)} \right),
\) donde $b^{(k)}$ es el sesgo y $\sigma(\cdot)$ la activación.  

Este mecanismo permite aprender filtros jerárquicos: los primeros sensibles a patrones locales (armónicos, periodicidades) y los posteriores a representaciones más abstractas que diferencian especies o condiciones de grabación.

\subsubsection{Aplicación y robustez en HumBugDB}

En HumBugDB, las CNNs se entrenan sobre ventanas de audio convertidas en mel-spectrogramas, que concentran la estructura espectral relevante en un dominio perceptual comprimido.  
Campos et al. (2022) mostraron que, incorporando \emph{data augmentation} acústico —mezcla con ruido de fondo, variaciones de tono y estiramiento temporal—, los modelos aprenden representaciones invariantes a perturbaciones comunes, mejorando su desempeño en ambientes ruidosos \cite{Campos2022}.  

Los \emph{kernels} se optimizan mediante retropropagación y no se fijan como filtros predefinidos, lo que permite a la red adaptarse a patrones específicos de mosquitos.  
La normalización por lotes y el \emph{dropout} introducen estocasticidad que previene el sobreajuste.  
Además, al mantener dropout activo en inferencia, se generan múltiples predicciones por entrada, lo que habilita la estimación de incertidumbre asociada a cada salida \cite{GalGhahramani2016}.  

\subsubsection{Resultados}

Los estudios iniciales reportaron que las CNNs superan ampliamente a clasificadores tradicionales (SVMs, Random Forests) cuando se entrenan con mel-spectrogramas de HumBugDB.  
Además de detectar la presencia de mosquitos, los modelos pueden diferenciar especies, lo cual aporta valor epidemiológico para programas de control de vectores \cite{HumBugDB, Campos2022}.  
La combinación de \emph{data augmentation}, kernels entrenables y regularización bayesiana mediante dropout convierte a estas arquitecturas en una herramienta robusta para la detección acústica en condiciones reales.
\insertarfigura{fig/cnn_humbugdb.png}
{Arquitectura típica de una CNN utilizada en HumBugDB. 
	El pipeline incluye: entrada de mel-spectrogramas, convoluciones con activación ReLU y max-pooling, 
	aplanamiento de características y capas densas para la clasificación. 
	Manteniendo dropout activo también en la fase de inferencia es posible estimar la incertidumbre asociada a cada predicción.}
{fig:cnn_humbugdb}{0.1}


\subsection{Relevancia y aplicaciones}

El uso de \textbf{CNNs} en HumBugDB responde a una necesidad concreta: 
\textbf{detectar y clasificar señales acústicas en entornos altamente ruidosos}.  
En este caso se trata de zumbidos de mosquitos, pero el problema es más general: 
identificar patrones espectro-temporales que se mantienen estables aun cuando el ambiente introduce 
variaciones de fondo, superposición de fuentes o degradaciones de calidad.

La relevancia del análisis es clara: si bien métodos clásicos como filtros adaptados o PLLs 
pueden detectar frecuencias dominantes, no ofrecen la flexibilidad de discriminar entre especies 
o de adaptarse a entornos cambiantes.  
Las CNNs, en cambio, aprenden \textbf{representaciones jerárquicas} que capturan tanto 
patrones locales (armónicos, periodicidad) como características más abstractas que permiten generalizar.  

El problema que se intenta resolver es: \textit{¿cómo construir un detector robusto y práctico 
	que funcione no sólo en condiciones de laboratorio, sino también en ambientes reales, 
	con ruido, voces, tráfico o viento?}.  
La respuesta está en combinar arquitecturas convolucionales con técnicas de regularización 
y estimación de incertidumbre (MC Dropout), lo que brinda un detector más confiable.

Más allá de los mosquitos, estas mismas técnicas pueden aplicarse a:  
\begin{itemize}
	\item \textbf{Detección de disparos y explosiones} en seguridad pública, para activar alertas tempranas en ciudades.  
	\item \textbf{Reconocimiento de llamadas de ballenas y delfines} en estudios oceánicos, incluso bajo el ruido de barcos.  
	\item \textbf{Identificación de defectos en soldaduras} a partir del sonido emitido durante el proceso industrial.  
	\item \textbf{Monitoreo de cultivos agrícolas} detectando plagas a través de sonidos de insectos distintos a los mosquitos.  
	\item \textbf{Diagnóstico de fallas en vehículos eléctricos}, distinguiendo ruidos anómalos de los motores o inversores.  
\end{itemize}

De esta forma, las CNNs no sólo representan un salto frente a técnicas clásicas en detección de mosquitos, 
sino que muestran cómo el \textbf{aprendizaje profundo} puede extenderse a un amplio rango de problemas 
donde las señales acústicas contienen información crítica que debe ser detectada en tiempo real.

%\subsection{CNNs aplicadas a HumBugDB}
% Archivo: cap_dl_2.tex
% Contenido: Uso de CNNs para clasificación de mosquitos en HumBugDB. 
% Técnicas de robustez frente al ruido (augmentaciones, kernels entrenables).
% Referencias: HumBugDB (2021), Campos (2022).


\subsection{Modelos secuenciales}
\label{sec:dl_seq_bayes}

\subsubsection*{RNN y LSTM}

Otra familia de arquitecturas fundamentales en audio son las \emph{Redes Neuronales Recurrentes} (RNNs), diseñadas para procesar secuencias temporales.  
En una RNN, la salida en cada instante $t$ depende no solo de la entrada $x_t$, sino también del estado oculto $h_{t-1}$, lo que permite modelar dependencias a lo largo del tiempo.  

El entrenamiento de RNNs estándar enfrenta el problema del \emph{gradiente desvanecido o explosivo}.  
Durante la retropropagación a través del tiempo (BPTT), los gradientes se multiplican repetidamente por derivadas de funciones de activación.  
Si estas son menores que uno, los gradientes tienden a cero (\emph{vanishing}); si son mayores que uno, crecen exponencialmente (\emph{exploding}).  
Como consecuencia, la red tiene dificultades para aprender dependencias de largo plazo de manera estable.

Para mitigar este problema se introdujeron variantes como las \emph{Long Short-Term Memory} (LSTM) y las \emph{Gated Recurrent Units} (GRU), que incorporan compuertas de control del flujo de información:
\begin{itemize}
	\item La \textbf{puerta de olvido} decide qué parte del estado anterior se descarta.  
	\item La \textbf{puerta de entrada} regula qué información nueva se incorpora a la memoria.  
	\item La \textbf{puerta de salida} determina qué fracción del estado interno influye en la salida actual.  
\end{itemize}

Gracias a este diseño, las LSTM mantienen gradientes estables en secuencias largas, permitiendo capturar dependencias temporales de distinto alcance \cite{DL_MusicReview}.  
Estas arquitecturas han sido ampliamente utilizadas en tareas de \emph{Music Information Retrieval} (MIR), como clasificación de género musical, transcripción automática o separación de fuentes, y continúan siendo un estándar en el modelado de series temporales \cite{DL_MusicReview}.


\subsection{Relevancia y diferencias frente a CNNs}

Tanto las CNNs como las RNNs trabajan con información temporal, pero de manera distinta.  
En un espectrograma, una CNN aprende patrones \textbf{locales} en bloques tiempo–frecuencia (armónicos, texturas, repeticiones), 
y puede combinarlos jerárquicamente para formar representaciones más complejas.  
Sin embargo, su ventana de contexto queda limitada por el tamaño de los filtros y la profundidad de la red.

Las RNNs, en cambio, procesan la secuencia de cuadros de manera explícita: el estado oculto transporta información entre instantes, 
lo que permite modelar dependencias \textbf{a más largo plazo} que las que cubre un kernel convolucional.  
Por eso resultan útiles cuando interesa capturar variaciones lentas o estructuras que se extienden en el tiempo 
(como la evolución del zumbido de un insecto o el fraseo en música).

En síntesis:
\begin{itemize}
	\item \textbf{CNNs}: detectan patrones locales robustos en el plano tiempo–frecuencia.  
	\item \textbf{RNNs/LSTM}: modelan cómo esos patrones se encadenan a lo largo de secuencias extensas.  
\end{itemize}

Hoy en día, ambos enfoques se consideran complementarios y a menudo se integran en arquitecturas híbridas.

% Archivo: cap_dl_3.tex
% Contenido: 
% - RNN/LSTM: captura de dependencias temporales en audio.
% - Dropout como inferencia bayesiana para incertidumbre.
% Referencias: Gal & Ghahramani (2016), Music DL Review (2017, 2023).



\subsection{Deep Learning para dirección de arribo y beamforming}
\label{sec:dl_doa}

La estimación de dirección de arribo (DoA) es un problema central en radar, comunicaciones inalámbricas y procesamiento acústico.  
Los algoritmos clásicos como MUSIC destacan por su alta resolución angular, pero presentan limitaciones conocidas: baja robustez con pocas muestras (\emph{snapshots}), sensibilidad al ruido y dificultad para localizar fuentes coherentes.  
Estas restricciones reducen su eficacia en escenarios prácticos.

Un enfoque reciente, denominado \textbf{DLMUSIC}, combina técnicas de \emph{Deep Learning} con el marco clásico de MUSIC \cite{BaiChen2023_DLMUSIC}.  
La propuesta introduce dos innovaciones principales:
\begin{itemize}
	\item \textbf{Estimación de la covarianza:} en lugar de usar la matriz de covarianza empírica, se entrenan redes convolucionales (CNNs) para predecir una versión más estable y robusta, menos sensible al ruido y al número de snapshots.  
	\item \textbf{Espectro objetivo suavizado:} en vez de aproximar el espectro MUSIC real (que puede presentar picos distorsionados y lóbulos espurios en baja SNR), se utiliza un \emph{espectro idealizado} compuesto por funciones Gaussianas centradas en las direcciones verdaderas de llegada.  
	Esta representación suavizada actúa como objetivo de entrenamiento, guiando a la red hacia picos definidos y reduciendo falsos máximos.
\end{itemize}

Tras la estimación de la covarianza por las CNNs, el algoritmo procede con la descomposición en autovalores y la separación de subespacios propia de MUSIC, preservando así la interpretabilidad del método clásico.  
De esta manera, DLMUSIC no reemplaza la teoría de subespacios, sino que corrige las etapas más frágiles mediante aprendizaje profundo.

De acuerdo con los resultados reportados, este enfoque supera al MUSIC clásico en escenarios de \textbf{baja SNR}, con \textbf{pocos snapshots} y en presencia de \textbf{fuentes coherentes} \cite{BaiChen2023_DLMUSIC}.  
Ello demuestra cómo el \emph{Deep Learning} puede complementar a los algoritmos tradicionales, reforzando su desempeño en condiciones adversas y acercándolos a un uso más confiable en aplicaciones modernas.



\subsection{Relevancia y aplicaciones}

La importancia de enfoques como \textbf{DLMUSIC} radica en que resuelven varias de las limitaciones prácticas 
de MUSIC clásico.  
En escenarios reales, las señales rara vez cumplen las condiciones ideales: 
se dispone de pocas muestras, el ruido no es despreciable y las fuentes pueden estar correlacionadas.  
En esos casos, MUSIC pierde resolución o incluso falla en detectar las direcciones correctas.  

El aporte de DLMUSIC es reemplazar las etapas más frágiles —estimación de covarianza y espectro de referencia— 
por versiones aprendidas y suavizadas mediante \emph{deep learning}.  
Esto no elimina la interpretabilidad del método subespacial, pero lo hace más \textbf{robusto frente a ruido, 
	coherencia y baja cantidad de datos}.  
En la práctica, significa pasar de un algoritmo que funciona bien solo en condiciones de laboratorio 
a una herramienta con proyección en campo abierto.

Más allá de los mosquitos, esta combinación de aprendizaje profundo con teoría de subespacios 
puede aplicarse en dominios menos tradicionales, por ejemplo:
\begin{itemize}
	\item \textbf{Ubicación de fugas de gas} en plantas industriales usando arreglos de micrófonos para detectar ultrasonidos.  
	\item \textbf{Localización de jugadores} en estadios mediante arreglos acústicos para análisis deportivo en tiempo real.  
	\item \textbf{Identificación de puntos calientes en líneas de transmisión eléctrica} combinando antenas de RF y sensores térmicos.  
	\item \textbf{Seguimiento de fauna submarina} (peces o crustáceos) mediante redes de hidrófonos de bajo costo.  
\end{itemize}

Así, DLMUSIC no sólo refuerza la robustez de MUSIC en condiciones adversas, 
sino que abre la puerta a su uso en aplicaciones emergentes donde antes la estimación de dirección 
no era viable por falta de SNR o de datos.

% Archivo: cap_dl_4.tex
% Contenido: Integración de redes neuronales con algoritmos clásicos (MUSIC, Capon).
% Uso en dirección de arribo y beamforming adaptativo.
% Referencias: Gupta (2015), Choi et al. (2017), Music DL Review (2023), videos Python for Everybody.
